% Inclusion fragment. Exact declaration ledger: grh_provenance.json.
% New elementary deductions are distinguished from retained arithmetic inputs.
\subsection{What changes for Dirichlet GRH}

\paragraph{GH1. Scope and fixed target.}
The exported general statement concerns every Dirichlet character, of every
positive modulus: $L(s,\chi)\ne0$ for $\Re s>1/2$, with the principal
character's pole at $s=1$ excluded.  It also asserts that every zero in
$0<\Re s<1$ lies on the critical line.  Unrestricted Dedekind or
automorphic GRH is not asserted by these declarations.

Fix a primitive character $\psi$ of conductor $M$ and order $o$.
Let $\psi_j$ be the primitive inducer of $\psi^j$, $0\le j<o$.
These $o$ characters are distinct: equal primitive inducers induce the
same character back to modulus $M$.  They are closed under inversion,
$\psi_0=\mathbf1$, and $\psi$ occurs.  Thus the source's
\texttt{Level.cyc} is exactly the fixed positive family
\[
 H_\psi(x)=G(x)+\frac1{3o+1}\sum_{j=1}^{o-1}\Re G_{\psi_j}(x).
 \tag{GH2}
\]
Here the centered sums $G_\chi$ are those already used in AN3.  No target
character changes as the auxiliary fields grow.  The arithmetic goal is
the same window bound as before, with constants allowed to depend on this
fixed $\psi$; AN3 already proves that it excludes every constituent zero
to the right of $1/2$.

\paragraph{GH3. An explicit host and the constant $c$.}
Write $M=2^aM_0$, $M_0$ odd, and use the Chinese remainder theorem to write
$\psi=\psi_2\psi_{\rm odd}$.  Put
\[
 K_\psi=\mathbb Q(\zeta_M)^{\ker\psi},\quad
 T_e=\mathbb Q(\zeta_{2^{e+2}})^+,\quad
 c=\operatorname{lcm}\bigl(\operatorname{ord}\psi_{\rm odd},
                  \operatorname{ord}(\psi_2(-1))\bigr).
\]
For $e\ge1$ and $e+2\ge a$, and distinct primes
$q\equiv3\pmod8$, $d\equiv7\pmod{16}$ with $\gcd(M,qd)=1$,
\[
 E_e=T_eK_\psi(\sqrt{-q},\sqrt{-d}),\qquad
 [K_\psi\cap T_e:\mathbb Q]=o/c,\quad
 [E_e:\mathbb Q]=4c2^e,\quad n=c2^e.                 \tag{GH3}
\]
Indeed the character group of $K_\psi$ is $\langle\psi\rangle$.
Its member $\psi^j$ belongs to the character group of $T_e$ precisely
when its odd-conductor component is trivial and its $2$-component is
even: $\operatorname{ord}(\psi_{\rm odd})\mid j$ and
$\operatorname{ord}(\psi_2(-1))\mid j$.  Their intersection therefore
has $o/c$ elements.  The two new quadratic characters are independent
of this base group by CRT at $q$ and $d$, giving the factor four.
In the ambient cyclotomic field, fixing $E_e$ means exactly
\[
 u\equiv\pm1\pmod{2^{e+2}},\quad \psi(u)=1,
 \quad J_q(u)=J_d(u)=1.
\]
This is \texttt{IsChiField}; it identifies the actual source field up to
$\mathbb Q$-isomorphism.  The source's auxiliary modulus is also explicit:
$Q_S=\operatorname{lcm}(2^{e+2},M)$, since its definition gives both
$2^{e+2}\mid Q_S$ and $Q_S\mid\operatorname{lcm}(2^{e+2},M)$,
while primitivity gives $M\mid Q_S$.
The proof explains the eventual constant in \texttt{ordAny} without
an unspecified stabilization stage.  In particular $1\le c\le o$.
For $M$ odd, $c=o$; for a purely $2$-primary even character, $c=1$.
Total evenness alone does not give $c=1$: $\psi=\chi_{-4}\chi_{-3}$
is even but has $c=2$.

\paragraph{GH4. The finite character calculation.}
The quotient characters used by the source are
$\psi^jJ_q^iJ_d^k$ ($0\le j<o$, $i,k\in\{0,1\}$).
To see completeness, a character trivial on the first source factor
descends to $(\mathbb Z/M\mathbb Z)^\times/\ker\psi$; every character
of this cyclic group is a power of $\psi$.  The remaining two sign
factors supply $i,k$.  On integers coprime to $2Mqd$,
\[
 \sum_{j=0}^{o-1}\sum_{i,k=0}^1
   (-1)^k\psi(n)^jJ_q(n)^iJ_d(n)^k
 =o\,\mathbf1_{\psi(n)=1}(1+J_q(n))(1-J_d(n)).       \tag{GH4}
\]
This follows by multiplying three finite geometric sums.
It is an exact support identity, not a square-root estimate.
At ramified primes the primitive representatives can have different
values from the characters extended by zero modulo $M$; the displayed
coprimality condition must be retained.
The omitted source characters have $(i,k)\ne(0,0)$, exactly the twists
controlled by \texttt{PairGood}.  Its bounds, and the cohomological
window theorem, remain arithmetic inputs here.

\paragraph{GH5. A fixed pair after the degree rescaling.}
Assume the full \texttt{PairEx} and \texttt{VWinAt} inputs for $\psi$.
Fix $0<\varepsilon<1/52$, put $\alpha=2+8\varepsilon$ and
$\kappa=c^{1/\alpha}$.  There are $J\ge1$ and arbitrarily large $Z$
with a \emph{single} $Y,q,d$ such that every $x\in[Z,Z^{3/2}]$ has
\[
 \begin{gathered}
 x^{1/J}\le Y\le x,\quad \operatorname{PairGood}(M,\varepsilon,Y,x;q,d),\\
 x^\alpha\le c2^{e(x/\kappa)}\le2x^\alpha,\qquad
 \mathcal V(E_{e(x/\kappa)},\lfloor x^{400}\rfloor)
       \le\exp(\tfrac12x^\varepsilon),
 \end{gathered}                                                    \tag{GH5}
\]
where $e(t)=\lceil\alpha\log_2t\rceil$ for $t\ge1$.
For a proof, obtain $J$ from the cohomological input, then select
$q,d$ good on $[Y,X]$ with $Y=X^{1/J}$.  That input supplies
$Y\le W$, $W^2\le X$ and its bound through degree $\lfloor t^{401}\rfloor$
for all $W\le t\le W^2$.  Choose $X$ large enough that
$W\ge\max(\kappa^3,\kappa^{400},1)$, and set $Z=\kappa W$.
Then $x\le Z^{3/2}\le W^2\le X$, $t=x/\kappa\in[W,W^2]$, and
$x^{400}=\kappa^{400}t^{400}\le t^{401}$.
Also $c\,t^\alpha=x^\alpha$ and $t^\varepsilon\le x^\varepsilon$.
These are all the required comparisons.  Any fixed lower bound on $e$
can be met by increasing $X$.  The quantifiers fix $q,d$ before $x$;
only the tower exponent varies.

\paragraph{GH6. Primitive to arbitrary characters, including the boundary.}
Let $\chi$ modulo $N$ be induced by a primitive $\chi^*$ of conductor
$f\mid N$.  Away from a principal pole,
\[
 L(s,\chi)=L(s,\chi^*)
       \prod_{\substack{p\mid N\\p\nmid f}}
             (1-\chi^*(p)p^{-s}).                         \tag{GH6}
\]
For $\Re s>0$, every displayed factor is nonzero, because
$|\chi^*(p)p^{-s}|=p^{-\Re s}<1$.
Thus AN3 applied to GH2 for every primitive target yields nonvanishing
for all characters to the right of $1/2$.  At $s=1$ a nonprincipal
$\chi$ has nonprincipal $\chi^*$, so the same identity applies and
the classical theorem \texttt{LFunction\_ne\_zero\_of\_one\_le\_re}
for nonprincipal characters covers this point; a principal character
has a pole and is excluded.  Within $0<\Re s<1$ the primitive functional
equation reflects a zero to $1-s$ for $(\chi^*)^{-1}$: the conductor,
root-number, and gamma factors there are nonzero.  A zero left of
$1/2$ would contradict right-half-plane nonvanishing.  Added Euler
factors can have zeros on $\Re s=0$, which is why the all-character
zero statement is about the open strip.

\paragraph{GH7. Shared work and the remaining arithmetic.}
The local polynomial, period, Gram, integral-index, slope, and positive-family
analytic arguments are shared.  Ordinary RH sets $M=o=c=1$ and
$H_\psi=G$.  Dirichlet GRH additionally requires the full fixed-$M$
pair-selection theorem, the host $K_\psi$, and cohomological windows
uniform across its growing tower.  These latter inputs do not follow
from the elementary character and degree calculations above; their
precise extra parameters are recorded in the companion report.
